%!TEX root = ../hutbthesis_main.tex

\chapter{总结与展望}

\section{总结}

本文就依靠FastMCP的人车仿真器交互控制系统开展设计与达成，把自然语言交互，轻量化部署，标准化接口当作核心目标，创建起一套适合自动驾驶仿真和虚拟场景编辑的完备系统，在研究期间，把MCP协议当作通信核心，连通了大模型，人车仿真器和虚幻引擎之间的数据链路，达成了从语音，文本指令到仿真执行的全过程闭合控制，系统整体采取分层架构设计，包含用户交互层，MCP服务层和仿真执行层，各个模块解耦分明，接口统一，有着较好的可守护性和扩展性。

系统具备诸多关键功能，其涉及车辆及无人机的动态控制，天气与环境参数化调节，多传感器数据采集，Unreal场景与蓝图程序化编辑，自然语言指令解析映射以及轻量化一键部署等方面，经由模块化工具封装并设计标准化接口之后，用户无需编写代码就能达成仿真场景创建，无人车控制，环境编辑等复杂操作，极大地减小了自动驾驶仿真平台的使用难度，从检测结果来看，该系统运行稳定，指令响应快速，场景适配性良好，可以满足教学，科研以及算法验证等多种场景的需求。

本文章把FastMCP协议同人车仿真融合，给自动驾驶仿真带来轻量化，智能化，低代码的新方案，该系统既改善了仿真操作的便捷性和效率，又为依靠自然语言的智能仿真交互赋予可复用的设计想法和工程操作方法，对于相关领域的研究和应用具备一定的参考意义。

\section{展望}

本系统已具备核心功能并完成验证，不过由于时间及软硬件条件所限，仍存有改进与拓展之处，日后可围绕功能加强，性能改善，场景扩充等方面加以完善，从而让系统变得更为实用，通用。其一，可以进一步充实仿真场景库和模型资源，增添更多车型，无人机种类，行人模型以及交通设施等，以此来优化场景的真实感与丰富程度，进而满足更为复杂的人车交互仿真需求。

可以利用多模态交互能力，并结合语音合成，手势识别，视觉反馈等技术来达成更为自然流畅的人机交互效果，而且，要改良指令分析以及意图识别相关的算法，从而加强系统对于模糊指令，复杂长指令以及多轮对话的应对能力，使得该系统更好地符合实际应用情况，还可以增添仿真数据的记录，回放，分析以及报表生成这些功能，给算法验证赋予定量评价方面的支撑。

\enlargethispage{3\baselineskip}
在系统架构上，将来能够达成分布式部署，做到多用户共同在线协作编辑及仿真，加强并发处理能力，可以将系统同数字孪生，实景建模技术相融合，做到虚实交融仿真，从而进一步优化检测结果的可靠性，还可以把系统移植到嵌入式平台或者边缘计算设备当中，做到轻量化，便携式的部署，拓展其在实验室，教学实训，户外展示等更多场景中的应用可能性。
